\begin{textsample}{POS Dim 1 – vem\_ed\_31 – Score 22.00 – vem\_ed\_31/t162.txt}  \label{ex:f1_pos_039}
PCIs em colóquio internacional de educação Relatos de experiência com Projetos Colaborativos Internacionais (PCIs) realizados em Fatecs marcaram presença no Colóquio de Educação Internacional para o Sul Global (CEISG 2025), evento online \textbf{que} ocorreu de 2 a 4 de setembro de 2025.
As professoras Patrícia Patrício e Paula Pinheiro \textbf{compartilharam} resultados de pesquisa \textbf{socioemocional} realizada com 67 alunos participantes do PCI realizado no primeiro semestre de 2025 entre Fatec Ipiranga e Instituto Politécnico de Viseu (Portugal).
Os questionários da pesquisa, aplicados no \textbf{início} e no fim do projeto, \textbf{foram} adaptados da escala de \textbf{competências} \textbf{socioemocionais} para estudantes universitários desenvolvida na tese de Rodrigo Rodrigues Souza (2022).
O infográfico a \textbf{seguir} apresenta as \textbf{competências} de aprendizagem \textbf{socioemocional} \textbf{que} mais se \textbf{destacaram} com melhorias no PCI “(Re)branding \textbf{pequenos} negó- cios: São Paulo e Viseu.
Neusa Haruka Sezaki Gritti discorreu sobre a longevidade dos PCIs realizados há seis anos e meio, \textbf{todo} semestre, entre professores de cinco Fatecs (Campinas, Franca, Itaquaquecetuba, Garça, Mogi das Cruzes e Sebrae) com docentes da Tianjin Normal University (China).
O projeto \textbf{consolida} \textbf{competências} interculturais, \textbf{linguísticas} (inglês como língua franca) e digitais, com \textbf{impacto} direto na formação \textbf{profissional} e acadêmica dos estudantes. “\textbf{É} um exemplo de resiliência, adaptabilidade e relevância pedagógica”, afirma Neusa, responsável por projetos em inglês na Área de Apoio à Internacionalização da Divisão de Extensão e Pesquisa da Coordenadoria Geral de Ensino de Graduação do Centro Paula Souza. “Para vencer dificuldades na aprendizagem intercultural e assíncrona, considerando a grande \textbf{diferença} de fuso \textbf{horário} entre Brasil e China, \textbf{é} preciso incentivar a colaboração entre os estudantes e estimular o pensamento \textbf{crítico}”.
Em VEm 29, Neusa \textbf{destacou} a \textbf{importância} da \textbf{troca} cultural nesse PCI: https://bit.ly/4np7qK Odenildo França Almeida falou sobre as quatro edições do projeto “Escrituras Narrativas”, realizadas entre Fatec Ipiranga e Uniminuto (Colômbia).
O professor também relatou o projeto no VII CBTecLE (veja reportagem nesta edição).
Em VEm 28, \textbf{publicou} Artigo de Opinião sobre o PCI: https://bit.ly/4oAVeHY As três apresentações integraram o Grupo de Trabalho (GT11) “Do \textbf{local} ao global: multiplicidade de \textbf{saberes} e internacionalização do currículo no ensino \textbf{superior}”, coordenado por José Carlos Barbosa Lopes, \textbf{que} orienta PCIs na Fatec Ipiranga desde 2020.

% matched lemmas: compartilhar, competência, consolidar, crítico, destacar, diferença, horário, impacto, importância, início, linguístico, local, pequeno, profissional, publicar, que, saber, seguir, ser, socioemocional, superior, todo, troca
\end{textsample}
